\documentclass[a4paper,12pt,italian]{article}
\usepackage[italian]{babel}
\usepackage{longtable}
\usepackage{graphicx}
%\usepackage{lastpage}
\usepackage{fancyhdr}

\usepackage{setspace}
%\onehalfspacing

\pagestyle{fancy}
\fancyhead{} %annulla head di default
\fancyfoot{} %annulla foot di default

\newcommand{\titolo}{Ottimizzazione di CSP con Branch and Bound}
\newcommand{\autore}{Dissegna Stefano, Geremia Mirco}

%testa e piede
\lhead{\leftmark}
%\rhead{}
\cfoot{\thepage}%pagina \thepage\ di \pageref{LastPage}}

\title{\titolo}
\author{\autore}

\begin{document}

\begin{titlepage}
\begin{center}
%\vspace*{1in}
\includegraphics[scale=0.3]{unipd-logo.png} 

\par
\vspace{0.5in}
{\Huge Universit\`a degli studi di Padova}
\par
\vspace{0.5in}
{\LARGE Dipartimento di Matematica Pura ed Applicata}
\par
\vspace{0.3in}
\par
\vspace{0.5in}
{\huge \titolo}
\par
\end{center}
\vspace{0.5in}
\begin{center}
{\Large Studenti: \autore}
\par
\vspace{0.3in}
\end{center}
\par
\vspace{0.5in}
\begin{center}
Anno Accademico 2009-2010
\end{center}

\end{titlepage}

\tableofcontents
\newpage

\section{Scopo del progetto}

Risoluzione di problemi di ottimizzazione espressi sotto forma di
problema di vincoli binari con variabili con domini finiti. Le
performance del risolutore sono state testate su problemi generati
casualmente. Il risolutore \`e stato valutato con e senza la
propagazione di vincoli utilizzando AC come criterio di consistenza
locale.

\section{Rappresentazione dei domini}

I domini sono rappresentati in modo esplicito tramite una lista di
valori. Possono quindi essere rappresentati solo domini finiti. I
valori appartenenti ad un dominio possono essere solamente numeri
interi. Non si tratta di una grossa limitazione in quanto \`e sempre
possibile associare univocamente ad ogni intero un valore. Sui domini
\`e possibile compiere principalmente quattro tipi di operazioni:

\begin{enumerate}
\item Ottenere l'elenco dei valori appartenenti al dominio.
\item Ottenere una copia del dominio, utilizzato dall'alrgoritmo di
  ricerca per effettuare il backtracking.
\item Rimuovere, in modo distruttivo, dal dominio i valori che non
  sono consistenti con un dato vincolo. La relazione di consistenza
  utilizzata \`e l'AC (arc consistency) direzionale.
\item Controllare se il domino \`e fallito oppure no: un dominio
  fallito \`e un dominio che non contiene valori.
\end{enumerate}

\section{Rappresentazione delle variabili}

Ad ogni variabile \`e associato un nome ed un dominio di valori che
pu\`o assumere . Una variabile a cui \`e stato assegnato un valore
viene rappresentata come una variabile il cui dominio \`e un
singoletto contenente il valore.

\section{Rappresentazione dei vincoli}

I vincoli binari sono rappresentati in modo esplicito. Ogni vincolo
\`e composto da due variabili e da un insieme di coppie di valori
accettati tra queste due variabili. L'insieme \`e implementato
tramite una tabella di hash per poter eseguire velocemente il test di
appartenenza. I vincoli supportano le seguenti
operazioni:

\begin{enumerate}
\item Trasposizione: restituisce un nuovo vincolo uguale all'originale
  ma con le variabili trasposte.
\item Appartenenza: controlla se, data una coppia di valori per le due variabili, la
  coppia appartiene oppure no al vincolo. Data la rappresentazione
  esplicita dei vincoli, il controllo avviene semplicemente
  controllando che la coppia sia presente nell'insieme delle coppie
  del vincolo.
\end{enumerate}

\section{Funzione obiettivo ed euristica}

Le funzioni obiettivo ed euristica possono essere passate come
parametri all'algoritmo di ricerca come sottoclassi dell'interfaccia
\textit{Problem.Evaluator} implementando il metodo \textit{eval} che,
data una lista di variabili, ritorna la valutazione. Data la
rappresentazione uniforme delle variabili assegnate e delle variabili
non assegnate, l'interfaccia \textit{Problem.Evaluator} va bene come
base sia per la funzione obiettivo (che operer\`a su variabili i cui
domini sono singoletti) sia per l'euristica (che operer\`a invece su
domini con uno o pi\`u elementi).

\subsection{Funzione massimizzata}

La funzione obiettivo massimizzata durante i test \`e la somma dei
valori dei domini. L'euristica massimizza la somma dei valori pi\`u
alti dei domini. Se tutte le variabili sono assegnate, l'euristica
coincide con la funzione obiettivo, quindi \`e stata implementata solo
una funzione.

\section{Rappresentazione di un problema}

Un problema \`e rappresentato come:

\begin{itemize}
\item Una lista di variabili.
\item Una lista di vincoli tra le variabili.
\item Una funzione obiettivo.
\item Una funzione euristica.
\end{itemize}

\subsection{Generazione di un problema}

Il generatore di problemi prende quattro parametri (oltre alle
funzioni obiettivo ed euristica) che utilizza per generare un problema
random:

\begin{enumerate}
\item Numero di variabili del problema. I nomi delle variabili sono un
  numero progressivo.
\item Cardinalit\`a del dominio. Gli elementi di ogni dominio saranno
  numeri interi tra 0 e la cardinalit\`a (esclusa).
\item Densit\`a dei vincoli, un numero compreso tra 0 e 1. Man mano che vengono generati i vincoli
  binari tra tutte le possibili coppie di variabili, un vincolo viene
  con probabilit\`a pari alla densit\`a. Viene generato un numero
  casuale tra 0 e 1, e se non supera la densit\`a allora il vincolo
  viene accettato.
\item Strettezza dei vincoli, un numero compreso tra 0 e 1. Per ogni
  vincolo vengono generate tutte le coppie possibili di valori, ed
  ogni coppia viene accettata con una probabilit\`a pari alla
  strettezza.
\end{enumerate}

\section{Algoritmo di propagazione}

Durante la ricerca viene mantenuta la consistenza sugli archi
(AC). L'algoritmo utilizzato per la propagazione dei vincoli \`e
\textit{AC-1}. L'implementazione opera modificando distruttivamente i
domini delle variabili, quindi durante la ricerca \`e necessario
mantenere delle copie per poterli ripristinare durante il
backtracking. \textit{AC-1} utilizza anche i vincoli trasposti, che
vengono generati prima dell'inizio della ricerca e memorizzati, in
modo da non doverli ricalcolare ogni volta che viene richiamato l'\textit{AC-1}.
Per operare utilizza il metodo \textit{revise} dei vincoli, che
richiama poi \textit{removeInconsistent} sui domini. Questi due metodi
utilizzano il valore di ritorno per indicare se un dominio \`e stato
modificato oppure no. Se nessun dominio \`e stato modificato dopo
avere effettuato il revise su tutt i vincoli (trasposti e non),
l'algoritmo si ferma in quanto il problema \`e diventato \textit{AC}.

\section{Ricerca}
\subsection{Ripristinare i domini}
\subsection{Scelta della variabile}
\subsection{Scelta del valore}

\section{Test}

Per controllare la correttezza del programma, per ogni problema random
generato \`e possibile stampare l'equivalente script per il programma
\textit{Minion}. Questo permette di confrontare la soluzione trovata
da \textit{Minion} con quella del nostro programma.

Per la generazione, \`e stato aggiunto il metodo \textit{toMinion}
alle classi \textit{Variable}, \textit{BinaryConstraint},
\textit{Problem} e all'interfaccia \textit{Problem.Evaluator}. Questo
metodo genera la rappresentazione dell'oggetto in un formato testuale
adatto ad uno script \textit{minion}.

\section{Benchmark}

\subsection{Al variare di ...}
%% grafico + commento testuale

\end{document}
